\section{Weighted Fourier grading}\label{sec:grading}

We next construct an operator whose spectrum is fixed by the dimensions
of polynomial spaces, while its matrix entries respond to an arbitrary
positive weight. The construction is the cube instance of the usual degree
filtration for discrete multivariate orthogonal polynomials, whose
coordinate multipliers satisfy a three-term recurrence; compare
Xu~\cite[Section~3]{Xu2004}. We give the finite-dimensional argument in
the present notation.


Write $\Omega=\{-1,1\}^n$ and $N=2^n$. We work in
$\mathcal H=\mathbb C^\Omega$ with the counting inner product
\[
 \ip{a}{b}=\sum_{x\in\Omega}\overline{a(x)}b(x).
\]
Thus the point masses $\delta_x$ form an orthonormal basis. For an operator
$B$, write $B_{xy}=\ip{\delta_x}{B\delta_y}$ and let
\[
 \norm B_{\HS}^2=\sum_{x,y\in\Omega}\abs{B_{xy}}^2
\]
be its Hilbert--Schmidt norm squared; $\norm B_{\op}$ denotes its operator
norm. The Hilbert--Schmidt norm is independent of the orthonormal basis.

For $S\subseteq[n]=\{1,\ldots,n\}$, set
$\chi_S(x)=\prod_{j\in S}x_j$. These real characters satisfy
\[
 \ip{\chi_S}{\chi_T}=N\mathbf1_{\{S=T\}},\qquad
 \chi_S\chi_T=\chi_{S\mathbin\triangle T}.
\]
Indeed, a nonconstant character sums to zero by pairing vertices that
differ in one of its coordinates. The displayed product rule therefore
gives orthogonality, and the $N$ characters form a basis of the
$N$-dimensional space $\mathcal H$. Let
\[
 V_k=\operatorname{span}\{\chi_S:\abs S\le k\},\qquad 0\le k\le n,
\]
and set $V_k=\mathcal H$ for $k\ge n$. Membership in $V_k$ means that a
function has Fourier degree at most $k$. The character multiplication rule
shows that $a\in V_r$ and $b\in V_s$ imply $\overline b a\in V_{r+s}$.

Fix any function $w:\Omega\to(0,\infty)$, and let $D$ be multiplication
by $\sqrt w$. The identity
\[
 \ip{Da}{Db}=\sum_{x\in\Omega}w(x)\overline{a(x)}b(x)
\]
identifies the weighted inner product with the counting inner product
after applying $D$. We can therefore orthogonalize polynomial degrees for the weight
$w$ while retaining the orthonormal point basis used to count edges.
Form the nested subspaces
\[
 K_k=DV_k,\qquad K_{-1}=\{0\},\qquad
 E_k=K_k\cap K_{k-1}^{\perp}\quad(0\le k\le n).
\]
Let $\Pi_k$ be the orthogonal projection onto $E_k$. The spaces $K_k$
are nested, and $K_k=K_{k-1}\oplus E_k$. Since $D$ is invertible,
$\dim K_k=\dim V_k=\sum_{j=0}^k\binom nj$ and $K_n=\mathcal H$.
Subtracting consecutive dimensions gives
\[
\mathcal H=\bigoplus_{k=0}^n E_k,\qquad
 \dim E_k=\binom nk.
\]
Choosing an orthonormal basis in each $E_k$ and decomposing the image of
each basis vector gives, for every operator $B$,
\begin{equation}\label{eq:block-parseval}
 \norm B_{\HS}^2
 =\sum_{r,s=0}^n\norm{\Pi_s B\Pi_r}_{\HS}^2.
\end{equation}
Define the grading operator and its centered version by
\begin{equation}\label{eq:grading}
 L=\sum_{k=0}^n k\Pi_k,\qquad M=L-\frac n2\Id.
\end{equation}
They are self-adjoint. Their prescribed eigenvalue multiplicities give
\begin{equation}\label{eq:total-energy}
 \norm M_{\HS}^2=\sum_{k=0}^n\binom nk\left(k-\frac n2\right)^2
 =\frac{nN}{4}.
\end{equation}
The last equality is the variance identity for a sum of $n$ independent
Bernoulli random variables of parameter $1/2$.

Let $\rho(x,y)$ denote Hamming distance. All pairs of vertices in what
follows are \emph{ordered}. When $w\equiv1$, the identity
$\sum_{y:\rho(x,y)=1}\chi_S(y)=(n-2\abs S)\chi_S(x)$ shows that $M$ is
minus one-half of the cube adjacency matrix. Thus every edge entry has
squared modulus $1/4$; the proof below controls the total deficit from
this value for arbitrary positive weights.

The key local property is that multiplying by a coordinate connects
only equal or adjacent grades.

For each $i\in[n]$, let $Z_i$ be multiplication by $x_i$. It is a
self-adjoint unitary and commutes with $D$. Multiplication by $x_i$ raises
Fourier degree by at most one, so $Z_iK_r\subseteq K_{r+1}$.
For $s>r+1$, the inclusion $K_{r+1}\subseteq K_{s-1}$ and the
orthogonality $E_s\perp K_{s-1}$ give $\Pi_sZ_i\Pi_r=0$.
Taking adjoints gives the opposite triangular vanishing, so
\begin{equation}\label{eq:tridiagonal}
 \Pi_sZ_i\Pi_r=0\qquad\text{whenever }\abs{s-r}>1.
\end{equation}

Together with the fixed total energy in \eqref{eq:total-energy}, this
restriction will control the entries of $M$ between vertices at Hamming
distance at least two.
